\subsection*{Limitations}\label{sec:limitations}
Despite \oursacro{} meeting our desiderata, some aspects remain unexplored.
Firstly, while we cover several common techniques for UE for molecules this is not an exhaustive list, e.g., Bayesian neural networks~\citep{lambBayesianGraphNeural2020} or normalizing flows~\citep{kohlerSmoothNormalizingFlows2021}.
Secondly, recent research \cite{fu2022forces} has shown a gap between force prediction and simulation performance, emphasizing the importance of stability for backbones. However, due to noise in active learning \cite{mittal2019parting, ren2021survey} and challenges in bridging this gap for many core architectures \cite{fu2022forces}, we leave a simulation benchmark for UE as future work.

\section{Conclusion}\label{sec:conclusion}
In this work, we approached the field of uncertainty estimation (UE) for molecular force fields by defining six key desiderata, three `physics-informed' and three `application-focused', that UE for force fields should satisfy.
In a survey, we analyzed previous works based on these desiderata and found that none of them satisfies all of them.
To fill the gap, we proposed \oursacro{}, a localized GP-based extension to existing GNNs.
In our experimental evaluation, we found \oursacro{} to fulfill all desiderata and provide significant empirical improvements over previous single-forward methods for UE such as Monte Carlo dropout or evidential regression.


\section*{Acknowledgements}
This project is supported by the German Federal Ministry of Education and Research (BMBF), grant no. 01IS18036B, and the Free State of Bavaria under the Excellence Strategy of the Federal Government and the Länder. This project is further supported by the DAAD programme Konrad Zuse Schools of Excellence in Artificial Intelligence, sponsored by the Federal Ministry of Education and Research.
